Much of the conversation about Generation Z and workplace ethics asks whether young employees care more about ethics than their predecessors, and the evidence gives that question little traction. This paper has asked instead where ethical complaints go once an employer is seen to break an ethical commitment. Drawing on psychological contract theory and the exit-voice-loyalty framework, it proposed that promises made in public are more likely to be felt as betrayed when broken and more likely to be contested in public, that a safe and credible internal channel limits this pull, and that exit and silence become the likely outcomes when neither venue seems open. It also set out two rival accounts of why this pattern might appear more often among the youngest employees, one based on cohort and one on career stage, and described designs that could separate them. The main takeaway is that venue deserves a place alongside values in research on ethics at work. The main limitation is that the model is untested and one of its central constructs, promise publicity, still lacks a measure. A useful next step would be a multi-wave study of newcomers of different ages, which could test the venue propositions and give a first answer to whether Generation Z is distinctive in this respect or simply new to the organization.
